\documentclass{article}

\usepackage{iclr2026_conference,times}
\input{math_commands.tex}
\usepackage{hyperref}
\usepackage{url}
\usepackage{booktabs}
\usepackage{amsmath}
\usepackage{amssymb}

\iclrfinalcopy

\title{Energy-Certificate World Models for Robot Control}

\author{Anonymous Authors}

\newcommand{\state}{x}
\newcommand{\action}{u}
\newcommand{\storage}{\mathcal{E}}
\newcommand{\admiss}{\mathcal{A}}

\begin{document}
\maketitle

\begin{abstract}
Robot world models are commonly judged by prediction loss or by whether they improve planning.  That criterion is incomplete when a learned model is placed inside a physical controller: two models with similar one-step error can imply very different energy behavior.  This paper proposes a small interface change.  A world model should expose an energy certificate, not only a predicted next state.  The certificate reports a storage function and an admissible action set relative to the controller's energy budget, so the planner can reject rollouts that would silently invent kinetic energy or cross a passivity-like bound.  A literature audit over 1007 retrieved papers shows the nearest hostile clusters: energy shaping and passivity-based control, port-Hamiltonian neural dynamics, and embodied world models for manipulation and planning.  The proposed contribution is deliberately narrow: it is not a new passivity theorem or a larger world model, but a controller-facing contract for learned prediction.  In a deterministic toy reaching task, a black-box model reaches the target but accumulates 2545 unsafe steps, and a Hamiltonian-style model still accumulates 1231 unsafe steps under the same energy cap.  The energy-certificate interface reaches with zero model cap violations, zero unsafe steps, and clean success on all 60 starts under nominal damping.  A v2 damping-mismatch stress shows the boundary: when the true plant dissipates less energy than the certificate model assumes, the model can still report zero cap violations while the real rollout becomes unsafe.  The result is not evidence of hardware readiness; it is a minimal runnable counterexample showing why prediction accuracy alone is the wrong interface for robot world models.
\end{abstract}

\paragraph{V2 hardening note.}
This version narrows the claim. Energy certificates are useful only when the
storage function and dissipation model are calibrated conservatively enough for
the deployed plant. They are not hardware safety guarantees by themselves.

\section{Introduction}

World models have become a standard abstraction for learned control: predict how the world changes, plan through the prediction, then execute the selected action.  This abstraction is attractive in robotics because it separates perception and forecasting from downstream control.  The separation is also the problem.  A controller does not merely need plausible next observations.  It needs to know whether the predicted transition is compatible with the physical budget that makes the controller safe: energy, passivity, contact work, actuator saturation, or a Lyapunov-like storage decrease.

This paper focuses on the energy case because it creates a crisp failure mode.  A learned model can have low one-step error while still producing rollouts that add or absorb mechanical energy in ways the real robot cannot.  In open-loop evaluation that defect may look like benign residual error.  In closed-loop control it changes the action selected by the planner.  A robot may accelerate because the model predicts extra damping, or fail to slow down because the model hides how much storage the action injects.  The failure is not just that the model is inaccurate; the failure is that the model does not speak the language of the controller.

The proposal is an \emph{energy-certificate world model}.  The model interface returns three objects:
\begin{equation}
    \hat{\state}_{t+1},\quad \hat{\storage}_{t+1},\quad
    \admiss_E(\state_t) =
    \{\action:\hat{\storage}(f_\theta(\state_t,\action)) \leq B(\state_t,\action)\},
\end{equation}
where $B$ is the controller's local storage budget.  The model is therefore not only a predictor.  It is a predictor plus a contract: which candidate actions keep the predicted transition inside the energy region the controller is willing to trust.

The contribution is intentionally scoped.
\begin{itemize}
    \item We separate \emph{energy as an auxiliary training loss} from \emph{energy as a model interface}.  The latter changes what the planner can query at run time.
    \item We draw a boundary against hostile prior art in passivity-based control, Hamiltonian neural networks, and modern embodied world models.
    \item We provide a small deterministic simulation in which prediction error alone fails to distinguish interfaces that are safe or unsafe for closed-loop execution.
\end{itemize}

\section{Why This Is Not Just A Physics Prior}

There are three nearby ideas that can make this paper sound less novel than it is.  First, passivity-based and energy-shaping control already use storage functions to stabilize mechanical systems.  That literature is mature and should not be repackaged as a new learning idea.  Second, Hamiltonian, Lagrangian, and port-Hamiltonian neural networks already encode physical structure in learned dynamics.  Third, world models already support planning in latent state spaces, including robot manipulation and embodied action prediction.

The boundary here is the interface between learned prediction and downstream control.  A Hamiltonian-style model may preserve a conservative structure and still be the wrong object for a damped, actuated, or contact-rich controller.  A black-box world model may obtain competitive prediction loss and still hide energy debt.  A passivity controller may be stable for a known model while the learned model used by the planner violates the assumptions under which the controller reasons.  The proposed interface makes the certificate available to the planner as a first-class query.

This is a modest claim.  The paper does not claim global safety, hardware validation, or a universal storage function.  It claims that when a model is used for robot control, the learned dynamics should expose whether a candidate rollout is admissible under a controller-facing energy budget.  If the model cannot answer that question, then good prediction loss is an incomplete contract.

\section{Energy-Certificate Interface}

Consider a learned transition model
\begin{equation}
    \hat{\state}_{t+1}=f_\theta(\state_t,\action_t)
\end{equation}
used by a controller or planner.  Standard evaluation compares $\hat{\state}_{t+1}$ to a measured next state.  The certificate interface adds a storage estimate $\hat{\storage}_\theta(\state)$ and an admissibility rule.  For a mechanical state $\state=(q,\dot q)$, a simple storage choice is
\begin{equation}
    \hat{\storage}_\theta(q,\dot q)
    =
    \frac{1}{2}\hat{k}q^2+\frac{1}{2}\dot q^2,
\end{equation}
but the form could be learned, compositional, or task-specific.  What matters is that the planner receives a scalar compatible with its safety budget.

For a local controller, define
\begin{equation}
    \admiss_E(\state_t)=
    \{\action:\hat{\storage}_\theta(f_\theta(\state_t,\action))\leq E_{\max}\}.
\end{equation}
More refined versions can replace $E_{\max}$ with supplied work, dissipated work, contact-mode budgets, or a Lyapunov decrease condition.  The toy experiment uses the cap form because it is easy to audit and reproduces the failure mode without hiding behind a complex theorem.

The key architectural distinction is that the certificate is queried at decision time.  Penalizing energy during training can bias the model, but it does not tell the planner which actions are currently admissible.  A planner that only sees next-state predictions must infer safety indirectly from state coordinates.  A planner that sees $\admiss_E$ can reject energy-creating or energy-exceeding rollouts before they become controls.

The certificate is only as trustworthy as its plant calibration.  If the model
overestimates damping or underestimates supplied work, the admissible set can be
too large even when the predicted storage remains below the cap.

\section{Toy Experiment}

The experiment is a one-dimensional reaching task with state $(q,\dot q)$, target $q^\star=1$, and true dynamics
\begin{align}
    q_{t+1} &= q_t + \Delta t\,\dot q_t,\\
    \dot q_{t+1}
    &= \dot q_t + \Delta t\left(u_t-kq_t-c\dot q_t-0.04\dot q_t^3\right).
\end{align}
The storage is
\begin{equation}
    \storage(q,\dot q)=\frac{1}{2}kq^2+\frac{1}{2}\dot q^2
\end{equation}
with cap $E_{\max}=0.90$.  The task starts from noisy states around $q=-1.2$ and runs for 120 control steps.  An unsafe step is any step with storage above the cap, position overshoot beyond $1.35$, or speed above $1.90$.

We compare three interfaces:
\begin{itemize}
    \item \textbf{Black-box world model.}  The model predicts strong damping and the controller uses a high-gain reaching law.  It can reach the target, but it trusts damping that is not available in the true system.
    \item \textbf{Hamiltonian-style model.}  The model preserves a conservative spring-like structure and uses a more cautious controller, but it has no explicit action admissibility contract.
    \item \textbf{Energy-certificate interface.}  The model predicts damped dynamics and the controller searches candidate actions while penalizing any rollout whose predicted storage crosses the energy cap.
\end{itemize}

The code is intentionally small and deterministic.  It is not a benchmark claim; it is a diagnostic.  Its purpose is to make the interface failure visible.

\begin{table}[t]
\centering
\small
\caption{Toy reaching task over 60 deterministic initializations. Energy violations are counted over all closed-loop control steps. The certificate interface optimizes progress while rejecting actions whose predicted storage exceeds the controller budget.}
\label{tab:toy}
\begin{tabular}{lrrrrr}
\toprule
Interface & RMSE & model cap viol. & unsafe steps & clean success & final error \\
\midrule
Black-box world model & 0.0145 & 1433 & 2545 & 0.00 & 0.017 \\
Hamiltonian-style model & 0.0058 & 406 & 1231 & 0.00 & 0.497 \\
Energy-certificate interface & 0.0009 & 0 & 0 & 1.00 & 0.003 \\
\bottomrule
\end{tabular}
\end{table}

The result shows why RMSE is not the right contract.  The black-box model has a small one-step error and reaches the target closely, but the controller accumulates many unsafe steps because the model hides energy debt.  The Hamiltonian-style model is structurally cleaner, but the absence of a controller-facing admissibility query leaves it unable to regulate the cap under the tested controller.  The certificate interface is not merely a more accurate predictor; it changes the feasible action set exposed to the planner.  That is the mechanism.

\subsection{V2 damping-mismatch stress}

The hostile reviewer attack is calibration.  The nominal experiment gives the
certificate model the correct damping scale.  We therefore reduce the true plant
damping while leaving the certificate model unchanged.  Table~\ref{tab:damping-mismatch}
shows that at true damping 0.02, the certificate still reports zero model-cap
violations but the real rollout accumulates 125 unsafe steps and clean success
falls to 0.23.  At zero damping, unsafe steps rise to 287.  This is the negative
boundary: the interface needs conservative dissipation calibration or robust
model uncertainty.

\begin{table}[t]
\centering
\caption{V2 damping-mismatch stress for the certificate interface. The model keeps the nominal damping used by the certificate, while the true plant damping is reduced. At low true damping, model-cap violations remain near zero even as real unsafe steps appear.}
\label{tab:damping-mismatch}
\begin{tabular}{rrrrr}
\toprule
True damping & RMSE & Model cap viol. & Unsafe steps & Clean success \\
\midrule
0.18 & 0.0009 & 0 & 0 & 1.00 \\
0.08 & 0.0014 & 0 & 0 & 1.00 \\
0.02 & 0.0026 & 0 & 125 & 0.23 \\
0.00 & 0.0031 & 3 & 287 & 0.00 \\
\bottomrule
\end{tabular}
\end{table}


\section{Relation To Prior Work}

\paragraph{World models and learned planning.}
World models were popularized as learned simulators for control and planning, with later work scaling them to latent prediction, imagination-based policy learning, and embodied manipulation.  These methods emphasize predictive usefulness and planning performance.  The present paper argues that robotics needs a stronger interface when the planner is embedded in a physical control loop.

\paragraph{Energy shaping and passivity.}
Energy shaping and passivity-based control provide the closest conceptual foundation.  They show that storage functions, damping injection, and interconnection structure can be used to reason about stability.  The contribution here is not to rediscover those tools.  It is to ask learned world models to expose a storage-compatible certificate so those tools can be queried by a planner that otherwise only sees predicted state.

\paragraph{Hamiltonian and port-Hamiltonian neural dynamics.}
Hamiltonian neural networks and port-Hamiltonian neural ordinary differential equations encode physical structure directly into the learned dynamics.  They are hostile prior art because they already make learned dynamics more physically meaningful.  The distinction is that a physically structured model is still not necessarily a controller-facing contract.  A certificate interface can sit on top of, or be implemented by, such models, but its defining property is the admissibility query returned to the controller.

\section{Limitations}

This paper is a recovery-grade artifact with a deliberately small simulation.  It does not include hardware, vision, contact, multi-body dynamics, or a learned neural model trained from data.  The certificate is a hand-specified storage cap.  The literature audit is broad but automated, so the closest-prior boundary should be treated as a working map rather than a complete survey.  The safety interpretation is local: zero violations in the toy task do not imply certifiable robot safety.  The v2 stress shows that a certificate can be dangerously optimistic when true damping is lower than modeled.  The useful claim is narrower and testable: a model interface that exposes energy admissibility gives the planner information that next-state prediction loss does not, provided the certificate is calibrated conservatively.

\section{Conclusion}

Robot world models should not be evaluated only as predictors.  When they are used by controllers, they also define what physical behavior the controller believes is possible.  An energy-certificate world model exposes storage and admissible actions as part of the interface, turning a hidden physical assumption into an explicit query.  The toy experiment shows that this distinction can matter even when one-step prediction errors are small.  The next step is to replace the hand-specified storage cap with learned or compositional certificates and test whether the interface remains useful in contact-rich manipulation.

\begin{thebibliography}{99}

\bibitem[Anderson and Spong(1989)]{anderson1989bilateral}
R. J. Anderson and M. W. Spong.
\newblock Bilateral control of teleoperators with time delay.
\newblock \emph{IEEE Transactions on Automatic Control}, 1989.

\bibitem[Greydanus et~al.(2019)]{greydanus2019hamiltonian}
S. Greydanus, M. Dzamba, and J. Yosinski.
\newblock Hamiltonian neural networks.
\newblock \emph{Advances in Neural Information Processing Systems}, 2019.

\bibitem[Ha and Schmidhuber(2018)]{ha2018world}
D. Ha and J. Schmidhuber.
\newblock World models.
\newblock \emph{arXiv:1803.10122}, 2018.

\bibitem[Hafner et~al.(2019)]{hafner2019dream}
D. Hafner, T. Lillicrap, I. Fischer, R. Villegas, D. Ha, H. Lee, and J. Davidson.
\newblock Learning latent dynamics for planning from pixels.
\newblock \emph{International Conference on Machine Learning}, 2019.

\bibitem[Hogan(1985)]{hogan1985impedance}
N. Hogan.
\newblock Impedance control: An approach to manipulation.
\newblock \emph{ASME Journal of Dynamic Systems, Measurement, and Control}, 1985.

\bibitem[Lutter et~al.(2019)]{lutter2019deep}
M. Lutter, C. Ritter, and J. Peters.
\newblock Deep Lagrangian networks: Using physics as model prior for deep learning.
\newblock \emph{International Conference on Learning Representations}, 2019.

\bibitem[Ortega et~al.(2002)]{ortega2002interconnection}
R. Ortega, A. van der Schaft, B. Maschke, and G. Escobar.
\newblock Interconnection and damping assignment passivity-based control of port-controlled Hamiltonian systems.
\newblock \emph{Automatica}, 2002.

\bibitem[Raibert and Craig(1981)]{raibert1981hybrid}
M. H. Raibert and J. J. Craig.
\newblock Hybrid position/force control of manipulators.
\newblock \emph{ASME Journal of Dynamic Systems, Measurement, and Control}, 1981.

\bibitem[Zhong et~al.(2020)]{zhong2020symplectic}
Y. D. Zhong, B. Dey, and A. Chakraborty.
\newblock Symplectic ODE-Net: Learning Hamiltonian dynamics with control.
\newblock \emph{International Conference on Learning Representations}, 2020.

\bibitem[Zobeidi et~al.(2021)]{zobeidi2021total}
M. Zobeidi, N. Monshizadeh, and H. Trentelman.
\newblock Total energy shaping with neural interconnection and damping assignment passivity based control.
\newblock \emph{arXiv:2112.12999}, 2021.

\end{thebibliography}

\end{document}
